\title{Paper Draft}
\author{Henry et al}
\date{} % keep if you want to suppress date

\maketitle

\section{Introduction}

% This was ripped straight from Shaposhnikov, but don't think I need it anyway:
% Aging is a complex biological process characterized by the progressive decline of physiological functions and increased susceptibility to disease, ultimately leading to death {Shaposhnikov, 2022 159}.
Over the past several decades, research across diverse model organisms, from yeast and nematodes to flies and mammals, has revealed that the rate of aging is not fixed but can be modulated by genetic, environmental, and pharmacological interventions \citep{gems2013}. 
Remarkably, many of the molecular pathways that regulate lifespan are highly conserved, suggesting that fundamental mechanisms of aging may be shared across taxa \citep{lopez-otin2013}.

While much of aging research has involved direct perturbation of these conserved pathways, an intriguing dimension of aging biology was revealed by the discovery that sensory perception can modulate lifespan independent of the physical acts associated with perceived stimuli \citep{pontillo2025, gendron2020, miller2020}. 
Studies in the nematode worm, \textit{Caenorhabditis elegans}, and the vinegar fly, \textit{Drosophila melanogaster}, have demonstrated that sensory cues related to food, mates, and even the presence of dead conspecifics can initiate physiological changes that influence patterns of aging \citep{gendron2014, chakraborty2019, libert2007,hernandez-lima2025, miller2022}. 
These effects are mediated by neuronal circuits that emanate from sensory tissues and interface with deeper regions of the central nervous system, ultimately engaging conserved neuromodulators, including biogenic amines and neuropeptides, that influence metabolism, stress resistance, and longevity \citep{chakraborty2019, miller2020, miller2022}. 

In at least some cases, sensory perception has been shown to modulate lifespan through the motivational states \textit{per se} that it influences. 
For example, the neuropeptide F system, homologous to mammalian neuropeptide Y, is a component of reward- and stress-related processing, and its activation can reduce starvation resistance and shorten lifespan \citep{ryvkin2024, gendron2014}.
Sensory perception of potential mates shortens lifespan in male flies through circuits involving NPF signaling, whereas successful mating reverses these effects \citep{gendron2014}, suggesting that an unresolved expectation of reward can influence lifespan. 
Insatiable hunger also slows aging, despite promoting behaviors normally associated with decreased lifespan\citep{weaver2023a}. 
These observations indicate that identifying the neural circuits that interpret and integrate sensory inputs, as well as how they generate conserved internal states that prescribe physiological outcomes, may provide new insight into healthy aging.

In \textit{Drosophila}, one neural substrate has recently emerged as a strong candidate for such an interpretive center: the ellipsoid body (EB). 
The EB resides within the central complex, a region that shares functional and organizational homology with the mammalian basal ganglia \citep{strausfeld2013}. 
The EB is subdivided into concentric domains innervated by six major classes of ring neurons (termed "ER" or "R" neurons) that process visual, thermal, mechanosensory, and circadian information \citep{hulse2021, buhl2021, yan2023}.
Exposure of flies to dead conspecifics, termed “death perception,” reduces fat stores, lowers starvation resistance, and shortens lifespan in a manner that requires visual function and serotonin signaling through 5-HT2A-expressing R2 and R4 neurons \citep{chakraborty2019}.
Insulin-like peptides, as well as the insulin-responsive transcription factor dFOXO, are required for aging effects of death perception, suggesting that EB ring neurons interface with conserved metabolic pathways implicated in lifespan control \citep{gendron2023}. 
Additional roles for EB ring neuron classes in locomotor control, sleep, visual memory, and state-dependent action selection, suggest that it is well-positioned to link myriad sensory representations to organismal state \citep{ofstad2011, yan2023, buhl2021, kottler2019}.

The extent to which the results in \textit{Drosophila} may illuminate similar processes in other taxa may hinge on the nature of the EB's role in influencing health-related outcomes. 
One possibility is that these neurons function primarily as passive conduits or gates for lifespan-altering sensory information: if their function is disrupted, flies may simply fail to perceive \citep{seelig2013, omoto2017} or internally represent salient stimuli, and phenotypes normally induced by those stimuli would not manifest. 
A more intriguing alternative is that the EB is an instructive modulator of aging-related processes, and direct manipulation of ring neuron activity, independent of external sensory input, would be sufficient to alter aging. 
It is this latter case that, if true, suggests translational potential for targeting central neuromodulatory and sensory-integration mechanisms in aging research.
% EH: You asked for a ref in this last sentence, but i'm not sure what you have in mind. An example where they targeted sensory integration mechanisms based on invert research, or?

Here, we use a variety of optogenetic tools, which enable temporally precise activation or inhibition of genetically-defined neuronal populations, to show that \textit{Drosophila} EB ring neuron populations are prescriptive modulators of aging that encode behaviorally meaningful information and orchestrate diverse transcriptional and metabolic responses throughout the animal. 
Our findings indicate that these neurons can act as active regulators of pro-longevity physiology rather than only as conduits for perception-dependent aging signals. 
Activation of R2 neurons, inhibition of R4d neurons, and simultaneous activation of R2 and R4 neurons each extends lifespan in males and females, but these manipulations produce divergent transcriptional responses in head tissues, including distinct effects on insulin signaling, stress-response, and metabolic pathways. 
Despite these molecular differences, we detect no consistent changes in feeding, fat storage, locomotor activity, or climbing ability that account for the longevity effects, arguing against a simple explanation based on overt energetic or behavioral shifts. 
Valence assays further show that R2 activation alone produces a negatively-valenced state, whereas combined R2/R4 activation produces a positively-valenced state, indicating that distinct neural states with different immediate behavioral consequences can converge on increased lifespan. 
Transcriptomic and metabolomic profiling of peripheral tissues further reveals that R2 activation drives systemic metabolic remodeling, including suppression of proliferative programs and upregulation of electrochemical and membrane homeostasis pathways not observed in head tissues, demonstrating that EB ring neuron activity induces tissue-specific responses throughout the fly.
Together, these findings support a model in which the ellipsoid body serves as an integrative neural locus through which sensory and state-related processes can influence aging via multiple downstream physiological programs.

